\documentclass[10pt]{article}
\usepackage[margin=1in]{geometry}
\usepackage{amsmath,amssymb}
\title{Disproof of Conjecture 3490: chromatic polynomials determine chromatic numbers}
\author{Local verified working submission}
\date{October 3, 2026}
\usepackage{url}
% Compact consistent title for a short mathematical note.
\makeatletter
\renewcommand{\maketitle}{\begin{center}
{\Large\bfseries\@title\par}\vspace{0.6em}
{\small\@author\par\@date}
\end{center}\vspace{0.5em}}
\makeatother
\setlength{\emergencystretch}{3em}
\begin{document}
\maketitle

\paragraph{Claim.}
Conjecture 3490 asserts the existence of graph families with identical
chromatic polynomial and spectral radius but different chromatic
numbers, with minimum order ten. We prove that even a single pair of
finite graphs with identical chromatic polynomial and different
chromatic numbers cannot exist. Thus the extra spectral-radius
condition and the stated minimum order cannot repair the assertion.

\paragraph{Definitions.}
For a finite graph $G$ and a nonnegative integer $k$, a proper
$k$-coloring is a function $V(G)\to\{0,\ldots,k-1\}$ assigning
different colors to adjacent vertices. The chromatic polynomial $P_G$
is characterized by
\[
 P_G(k)=\#\{\text{proper $k$-colorings of }G\}
 \qquad(k\in\mathbb N).
\]
The chromatic number $\chi(G)$ is the least $k$ admitting a proper
$k$-coloring. These conventions also cover the empty graph, whose
chromatic number is zero.

\paragraph{Proof.}
For every $k\in\mathbb N$,
\[
 G\text{ is $k$-colorable}\quad\Longleftrightarrow\quad P_G(k)>0.
\]
Consequently $P_G=P_H$ implies that $G$ and $H$ admit proper
$k$-colorings for exactly the same values of $k$. Their least such
values agree, so $\chi(G)=\chi(H)$.
Equivalently, if $\chi(G)<\chi(H)$, evaluation at $k=\chi(G)$
would give $P_G(k)>0=P_H(k)$, a contradiction; the reverse inequality
is excluded in the same way. No pair claimed in the conjecture exists.

\paragraph{Lean verification and semantic correspondence.}
The self-contained file \texttt{Main.lean} imports only \texttt{Std}.
Vertices are \texttt{Fin n}; the edge relation is arbitrary and
decidable, so the result applies in particular to every finite simple
graph. The list \texttt{words k n} enumerates the length-$n$ lists with
entries below $k$. Theorems \texttt{mem\_words} and
\texttt{words\_nodup} establish completeness and absence of
duplicates. The predicate \texttt{Proper} checks the edge constraints,
and \texttt{chromaticEvaluation} counts the proper assignments.
The theorem \path{evaluation_pos_iff_colorable} proves that this
count is positive exactly when a proper coloring exists.

\texttt{ChromaticNumberIs} uses the actual existence and minimality
definition. The theorem
\path{equal_evaluations_equal_chromatic_numbers} proves equality
of chromatic numbers for any two graphs with identical natural
coloring counts. Integer polynomials are represented explicitly by
coefficient lists and evaluated by Horner's rule in
\texttt{polynomialEval}. The predicate \texttt{IsChromaticPolynomial}
is precisely the counting characterization above. The final theorem
\path{not_polynomial_separating_pair_claim} refutes the
existence of two graphs sharing an actual integer polynomial with
that characterization and having distinct chromatic numbers.
Thus the polynomial-to-counting bridge is included in the formal
statement, and not left as an assumed numerical coincidence.

\paragraph{Reproduction and trust.}
Run \texttt{lean Main.lean} using Lean 4.19.0. The command exits
successfully. The principal results depend only on the standard Lean
axioms \texttt{propext}, \texttt{Classical.choice}, and
\texttt{Quot.sound}. There are no added axioms, placeholders, or
uses of \texttt{native\_decide}. The proof is uniform in the numbers
of vertices and colors and is not limited to a bounded search.
\end{document}
